% Options for packages loaded elsewhere
\PassOptionsToPackage{unicode}{hyperref}
\PassOptionsToPackage{hyphens}{url}
\PassOptionsToPackage{dvipsnames,svgnames,x11names}{xcolor}
\documentclass[
  italian,
  11pt,
]{article}
\usepackage{xcolor}
\usepackage[top=3cm, bottom=3cm, left=2.5cm, right=2.5cm]{geometry}
\usepackage{amsmath,amssymb}
\setcounter{secnumdepth}{5}
\usepackage{iftex}
\ifPDFTeX
  \usepackage[T1]{fontenc}
  \usepackage[utf8]{inputenc}
  \usepackage{textcomp} % provide euro and other symbols
\else % if luatex or xetex
  \usepackage{unicode-math} % this also loads fontspec
  \defaultfontfeatures{Scale=MatchLowercase}
  \defaultfontfeatures[\rmfamily]{Ligatures=TeX,Scale=1}
\fi
\usepackage{lmodern}
\ifPDFTeX\else
  % xetex/luatex font selection
\fi
% Use upquote if available, for straight quotes in verbatim environments
\IfFileExists{upquote.sty}{\usepackage{upquote}}{}
\IfFileExists{microtype.sty}{% use microtype if available
  \usepackage[]{microtype}
  \UseMicrotypeSet[protrusion]{basicmath} % disable protrusion for tt fonts
}{}
\usepackage{setspace}
\makeatletter
\@ifundefined{KOMAClassName}{% if non-KOMA class
  \IfFileExists{parskip.sty}{%
    \usepackage{parskip}
  }{% else
    \setlength{\parindent}{0pt}
    \setlength{\parskip}{6pt plus 2pt minus 1pt}}
}{% if KOMA class
  \KOMAoptions{parskip=half}}
\makeatother
\usepackage{longtable,booktabs,array}
\newcounter{none} % for unnumbered tables
\usepackage{calc} % for calculating minipage widths
% Correct order of tables after \paragraph or \subparagraph
\usepackage{etoolbox}
\makeatletter
\patchcmd\longtable{\par}{\if@noskipsec\mbox{}\fi\par}{}{}
\makeatother
% Allow footnotes in longtable head/foot
\IfFileExists{footnotehyper.sty}{\usepackage{footnotehyper}}{\usepackage{footnote}}
\makesavenoteenv{longtable}
\ifLuaTeX
\usepackage[bidi=basic,shorthands=off]{babel}
\else
\usepackage[bidi=default,shorthands=off]{babel}
\fi
\ifLuaTeX
  \usepackage{selnolig} % disable illegal ligatures
\fi
\setlength{\emergencystretch}{3em} % prevent overfull lines
\providecommand{\tightlist}{%
  \setlength{\itemsep}{0pt}\setlength{\parskip}{0pt}}
\usepackage{mathpazo}
\usepackage[T1]{fontenc}
\usepackage{mathtools}
\usepackage{bm}
\usepackage{booktabs}
\usepackage{longtable}
\usepackage{array}
\usepackage{multirow}
\usepackage{fancyhdr}
\pagestyle{fancy}
\fancyhead[L]{}
\fancyhead[R]{\leftmark}
\fancyfoot[C]{\thepage}
\usepackage[font=small,labelfont=bf]{caption}
\usepackage{framed}
\definecolor{shadecolor}{RGB}{248,248,248}
\usepackage{float}
\let\origfigure\figure
\let\endorigfigure\endfigure
\renewenvironment{figure}[1][2] {\expandafter\origfigure\expandafter[H]} {\endorigfigure}
\usepackage{etoolbox}
\AtBeginEnvironment{longtable}{\small}
\usepackage{bookmark}
\IfFileExists{xurl.sty}{\usepackage{xurl}}{} % add URL line breaks if available
\urlstyle{same}
\hypersetup{
  pdftitle={Machine Learning Notes},
  pdfauthor={Nicola Moscufo},
  pdflang={it},
  colorlinks=true,
  linkcolor={NavyBlue},
  filecolor={Maroon},
  citecolor={Blue},
  urlcolor={NavyBlue},
  pdfcreator={LaTeX via pandoc}}

\title{Machine Learning Notes}
\author{Nicola Moscufo}
\date{4 Marzo 2026}

\begin{document}
\maketitle

{
\hypersetup{linkcolor=Black}
\setcounter{tocdepth}{3}
\tableofcontents
}
\setstretch{1.25}
\section{Multilingual Natural Language Processing: Words and
Tokens}\label{multilingual-natural-language-processing-words-and-tokens}

\subsection{The Ontology of Words}\label{the-ontology-of-words}

In natural language processing (NLP), the distinction between word
instances and word types is fundamental to vocabulary management and
statistical modeling.

\subsubsection{Word Types vs.~Word
Instances}\label{word-types-vs.-word-instances}

\begin{itemize}
\tightlist
\item
  \textbf{Word Types (\(|V|\)):} These represent the set of distinct
  linguistic units (vocabulary) present in a corpus. If a corpus
  contains multiple instances of the word ``the,'' it contributes only
  once to the count of types. The number of unique types defines the
  \textbf{vocabulary size}, denoted as \(|V|\).
\item
  \textbf{Word Instances (\(N\)):} These represent the total count of
  running words in a text, where every occurrence is counted, regardless
  of whether the word has appeared previously. \(N\) corresponds to the
  total length of the corpus.
\end{itemize}

\subsubsection{Illustrative Example}\label{illustrative-example}

Consider the sentence: \emph{``They picnicked by the pool, then lay back
on the grass and looked at the stars.''}

\begin{itemize}
\tightlist
\item
  \textbf{Instances (\(N\)):} 16 (if punctuation is excluded).
\item
  \textbf{Types (\(|V|\)):} 14.

  \begin{itemize}
  \tightlist
  \item
    \emph{Note:} The word ``the'' appears three times (instances), but
    only once as a type.
  \end{itemize}
\item
  \textbf{Mathematical Relationship:} In any corpus, the relationship is
  defined by \(N \geq |V|\).
\end{itemize}

\subsection{Corpora and Language
Diversity}\label{corpora-and-language-diversity}

A \textbf{corpus} (plural: corpora) is a structured collection of text,
typically curated to represent a specific linguistic domain, register,
or speaker intent.

\subsubsection{The Multilingual
Challenge}\label{the-multilingual-challenge}

While there are approximately 7,000 distinct languages spoken globally,
contemporary NLP research suffers from significant ``Anglocentrism''---a
historical bias where algorithms are predominantly designed, trained,
and evaluated using English datasets. This leads to several technical
challenges:

\begin{itemize}
\tightlist
\item
  \textbf{Resource Inequality:} High-resource languages (e.g., English,
  Mandarin, Spanish) have massive amounts of digitized text, whereas
  low-resource languages lack the data necessary for robust statistical
  model training.
\item
  \textbf{Code-Switching:} This phenomenon occurs when a speaker or
  writer alternates between two or more languages or dialects within a
  single utterance or conversation. Standard monolingual tokenizers and
  language models struggle with the syntactic and semantic shifts
  inherent in code-switched text.
\end{itemize}

\subsubsection{Key Takeaways}\label{key-takeaways}

\begin{itemize}
\tightlist
\item
  The vocabulary size \(|V|\) is a property of the unique lexicon, while
  \(N\) is the measure of total corpus length.
\item
  NLP systems should ideally be evaluated across diverse linguistic
  contexts to avoid overfitting to the idiosyncrasies of English.
\item
  Code-switching necessitates flexible NLP pipelines capable of handling
  multiple linguistic structures simultaneously.
\end{itemize}

\newpage

\subsection{Tokenization}\label{tokenization}

Tokenization is the computational process of segmenting a continuous
stream of raw input text into discrete units known as \textbf{tokens}.
These tokens serve as the fundamental input units for subsequent NLP
tasks, such as vectorization, syntactic parsing, and machine
translation.

\subsubsection{Objectives of
Tokenization}\label{objectives-of-tokenization}

The primary goal is to map unstructured character sequences into
meaningful discrete units that capture enough semantic information while
maintaining a computationally manageable vocabulary size.

\subsubsection{Challenges in Multilingual
Tokenization}\label{challenges-in-multilingual-tokenization}

Tokenization strategies vary significantly across languages due to
differences in orthography and morphology:

\begin{itemize}
\tightlist
\item
  \textbf{White-space Tokenization:} Simple splitting by spaces (common
  in English). This approach fails in languages like Chinese, Japanese,
  or Thai, which do not use spaces as word boundaries.
\item
  \textbf{Morphological Complexity:} In agglutinative languages (e.g.,
  Turkish, Finnish), a single word can be composed of many concatenated
  morphemes. Tokenizing at the word level leads to an excessively large
  \(|V|\) and the ``out-of-vocabulary'' (OOV) problem.
\item
  \textbf{Subword Tokenization:} Modern approaches (e.g., Byte Pair
  Encoding (BPE), WordPiece) decompose words into smaller sub-units.
  This balances the representation of rare words while keeping the total
  vocabulary size \(|V|\) restricted, effectively mitigating OOV issues.
\end{itemize}

\subsubsection{Summary of Tokenization
Approaches}\label{summary-of-tokenization-approaches}

{\def\LTcaptype{none} % do not increment counter
\begin{longtable}[]{@{}
  >{\raggedright\arraybackslash}p{(\linewidth - 4\tabcolsep) * \real{0.3333}}
  >{\raggedright\arraybackslash}p{(\linewidth - 4\tabcolsep) * \real{0.3333}}
  >{\raggedright\arraybackslash}p{(\linewidth - 4\tabcolsep) * \real{0.3333}}@{}}
\toprule\noalign{}
\begin{minipage}[b]{\linewidth}\raggedright
Method
\end{minipage} & \begin{minipage}[b]{\linewidth}\raggedright
Mechanism
\end{minipage} & \begin{minipage}[b]{\linewidth}\raggedright
Primary Use Case
\end{minipage} \\
\midrule\noalign{}
\endhead
\bottomrule\noalign{}
\endlastfoot
\textbf{Word-based} & Split by whitespace/punctuation & Simple English
pipelines \\
\textbf{Character-based} & Split every character & Models with very
small \(|V|\) \\
\textbf{Subword-based} & Split by frequent n-grams &
Multilingual/General-purpose \\
\end{longtable}
}

\subsubsection{Key Takeaways}\label{key-takeaways-1}

\begin{itemize}
\tightlist
\item
  Tokenization is the bridge between raw character streams and numerical
  feature spaces.
\item
  The optimal tokenization strategy is language-dependent and must
  account for scripts without word delimiters and morphologically rich
  languages.
\item
  Subword tokenization is the industry standard for modern deep learning
  models, as it handles rare words and multilingual input more
  gracefully than pure word-level tokenization.
\end{itemize}

\begin{center}\rule{0.5\linewidth}{0.5pt}\end{center}

\section{NLTK and Linguistic
Morphology}\label{nltk-and-linguistic-morphology}

\subsection{NLTK and Word
Tokenization}\label{nltk-and-word-tokenization}

The Natural Language Toolkit (NLTK) provides standard utilities for
tokenization, often relying on the concept of \textbf{orthographic
words}.

\begin{itemize}
\tightlist
\item
  \textbf{Definition:} Orthographic words are defined by the conventions
  of a specific writing system. In many Western languages, the
  whitespace character acts as the primary delimiter.
\item
  \textbf{Limitations:} This approach is insufficient for languages that
  lack spaces (e.g., Chinese) or those that use complex cliticization
  (e.g., French ``l'avion'' or English contractions). Defining a
  ``word'' is therefore a non-trivial linguistic and computational
  decision.
\end{itemize}

\subsection{Vocabulary Growth (Heaps'
Law)}\label{vocabulary-growth-heaps-law}

A central question in computational linguistics is: \emph{How many word
types should be included in a vocabulary?}

Empirical observation shows that as the number of word instances (\(N\))
grows, the number of distinct word types (\(|V|\)) also grows. This
relationship is formalized by \textbf{Heaps' Law}:

\[|V| = kN^{\beta}\]

\begin{itemize}
\tightlist
\item
  \(|V|\): The vocabulary size (number of unique types).
\item
  \(N\): The number of running word instances.
\item
  \(k, \beta\): Free parameters determined empirically (typically
  \(0.5 < \beta < 0.8\) for natural language).
\item
  \textbf{Implication:} The vocabulary is an ``open'' set; increasing
  the corpus size indefinitely will continue to yield new, previously
  unseen word types (e.g., through compounding or neologisms).
\end{itemize}

\subsection{Morphology: The Study of
Morphemes}\label{morphology-the-study-of-morphemes}

Words are not necessarily atomic units; they are often composed of
smaller, meaningful building blocks called \textbf{morphemes}.

\begin{itemize}
\tightlist
\item
  \textbf{Root:} The primary morpheme conveying the core semantic
  content of a word (e.g., ``work'').
\item
  \textbf{Affix:} A morpheme attached to a root to modify its
  grammatical function or meaning (e.g., ``-ed'').
\end{itemize}

\subsubsection{Classifying Affixes}\label{classifying-affixes}

Affixes are generally categorized into two functional classes:

\begin{enumerate}
\def\labelenumi{\arabic{enumi}.}
\tightlist
\item
  \textbf{Inflectional Morphemes:}

  \begin{itemize}
  \tightlist
  \item
    \textbf{Function:} Express grammatical relationships (e.g., tense,
    number, aspect).
  \item
    \textbf{Properties:} Highly productive and regular. They do not
    typically change the word's part-of-speech category.
  \item
    \textbf{Examples:} Plural ``-s'', past tense ``-ed''.
  \end{itemize}
\item
  \textbf{Derivational Morphemes:}

  \begin{itemize}
  \tightlist
  \item
    \textbf{Function:} Create new lexemes, often changing the word's
    syntactic category (part-of-speech).
  \item
    \textbf{Properties:} Less predictable and often idiosyncratic. They
    may only attach to specific subsets of words (e.g., ``careful'' is
    valid, but ``sleepful'' is not).
  \item
    \textbf{Example:} ``Care'' (noun) \(\rightarrow\) ``Careful''
    (adjective) \(\rightarrow\) ``Carefully'' (adverb).
  \end{itemize}
\end{enumerate}

\subsection{Clitics}\label{clitics}

A \textbf{clitic} represents a middle ground between an independent word
and a bound morpheme.

\begin{itemize}
\tightlist
\item
  \textbf{Characteristics:} They function syntactically like independent
  words (representing full semantic concepts) but are phonologically
  bound to a ``host'' word.
\item
  \textbf{Examples:}

  \begin{itemize}
  \tightlist
  \item
    English: The contraction \emph{'ve} in \emph{``I've''} functions as
    the verb ``have'' but cannot stand alone.
  \item
    Romance Languages: In Spanish or Italian, object pronouns often
    attach to the end of verbs (e.g., \emph{``dámelo''} - ``give it to
    me'').
  \end{itemize}
\item
  \textbf{NLP Significance:} Tokenization algorithms must decide whether
  to treat clitics as part of the host (potentially bloating the
  vocabulary with variants like ``do'' and ``don't'') or to split them
  into separate tokens to simplify downstream syntactic analysis.
\end{itemize}

\subsubsection{Key Takeaways}\label{key-takeaways-2}

\begin{itemize}
\tightlist
\item
  \textbf{Heaps' Law} illustrates that vocabulary size is an open-ended
  function of corpus size, making fixed-vocabulary models prone to OOV
  errors.
\item
  \textbf{Morphemes} are the atomic units of meaning; distinguishing
  between inflectional and derivational morphology is essential for
  stemming and lemmatization tasks.
\item
  \textbf{Clitics} demonstrate the tension between orthographic
  representation and syntactic reality, requiring careful pre-processing
  in sophisticated NLP pipelines.
\end{itemize}

\newpage

\section{Subword Tokenization: Byte Pair Encoding
(BPE)}\label{subword-tokenization-byte-pair-encoding-bpe}

Modern NLP architectures (e.g., Transformers) rely heavily on subword
tokenization to bridge the gap between character-level and word-level
representations. \textbf{Byte Pair Encoding (BPE)} is the standard
algorithm for this task, as it effectively manages the vocabulary size
\(|V|\) while ensuring that no word is entirely ``out-of-vocabulary''
(OOV) by decomposing unknown words into known subword units.

\subsection{The BPE Mechanism}\label{the-bpe-mechanism}

BPE iteratively merges the most frequent adjacent pair of tokens. It
operates on a frequency-based greedy heuristic to build a vocabulary of
fixed size \(V\).

\subsubsection{1. Training Phase}\label{training-phase}

The training process induces a subword vocabulary from a training
corpus.

\begin{itemize}
\tightlist
\item
  \textbf{Initialization:} Start with a base vocabulary consisting of
  every unique character in the corpus (the ``character unigram'' set).
\item
  \textbf{Iteration:}

  \begin{enumerate}
  \def\labelenumi{\arabic{enumi}.}
  \tightlist
  \item
    Count the frequency of all adjacent pairs of symbols in the training
    corpus.
  \item
    Identify the most frequent pair \((A, B)\).
  \item
    Create a new token \(AB\) by merging \(A\) and \(B\).
  \item
    Add \(AB\) to the vocabulary.
  \item
    Replace all occurrences of \((A, B)\) in the corpus with \(AB\).
  \item
    Repeat until a predefined vocabulary size \(|V|\) is reached or a
    maximum number of merge operations is performed.
  \end{enumerate}
\end{itemize}

\subsubsection{2. Encoding Phase}\label{encoding-phase}

Once the vocabulary is induced, encoding new (unseen) text involves
applying the learned merge rules in the exact order they were created
during training.

\begin{itemize}
\tightlist
\item
  For any word, segment it into its constituent character units.
\item
  Iteratively apply the merge rules to the sequence.
\item
  If a word contains characters or sequences not encountered during
  training, the process eventually results in individual characters,
  ensuring the model never faces a completely unknown word.
\end{itemize}

\subsection{Structural Constraints}\label{structural-constraints}

BPE as traditionally implemented relies on specific assumptions about
text structure:

\begin{itemize}
\tightlist
\item
  \textbf{Boundary Restriction:} By default, BPE merges are constrained
  within word boundaries. The corpus is typically pre-tokenized by
  whitespace and punctuation before the BPE training begins to ensure
  that the algorithm does not merge across distinct words.
\item
  \textbf{Orthographic Dependency:} The reliance on whitespace
  pre-segmentation makes BPE highly dependent on the orthographic
  conventions of the source language.
\end{itemize}

\subsubsection{Limitations and
Considerations}\label{limitations-and-considerations}

\begin{itemize}
\tightlist
\item
  \textbf{Language Bias:} The ``whitespace-pre-tokenization'' assumption
  is theoretically problematic for languages like Chinese, Japanese, or
  Thai, where words are not delimited by spaces. Applying BPE to these
  languages requires either character-level pre-segmentation or
  specialized segmentation tools (e.g., morphological analyzers) before
  subword induction.
\item
  \textbf{Efficiency:} The number of merge operations directly impacts
  the sequence length. Too few merges result in character-level
  representations (long sequences, high computational cost), while too
  many result in word-level representations (high risk of OOV/sparsity).
\item
  \textbf{Deterministic Output:} Once the merge rules are fixed, the
  encoding process is deterministic, providing a stable input
  representation for deep learning models.
\end{itemize}

\subsubsection{Key Takeaways}\label{key-takeaways-3}

\begin{itemize}
\tightlist
\item
  \textbf{BPE} optimizes the vocabulary size by merging the most
  frequent character sequences iteratively, allowing the model to
  represent rare words as combinations of frequent subwords.
\item
  \textbf{Training} involves building a vocabulary by merging, while
  \textbf{Encoding} involves applying these rules to segment new text.
\item
  The \textbf{pre-tokenization} step (e.g., whitespace splitting) is a
  critical bottleneck for non-Western languages, necessitating
  adaptations for effective multilingual application.
\item
  Subword tokenization effectively balances the trade-off between the
  low-level granularity of characters and the high-level semantics of
  full words.
\end{itemize}

\newpage

\section{Managing Spelling Errors and String
Similarity}\label{managing-spelling-errors-and-string-similarity}

In NLP, handling Out-of-Vocabulary (OOV) terms---such as misspellings,
neologisms, or morphological variants---is critical for robust
performance. While subword tokenization (e.g., BPE) implicitly handles
OOV words by decomposing them into sub-units, this process can obscure
the original semantic intent. Traditional methods utilize string
distance metrics to map noisy input to valid vocabulary terms.

\subsection{Minimum Edit Distance (Levenshtein
Distance)}\label{minimum-edit-distance-levenshtein-distance}

The \textbf{Minimum Edit Distance (MED)} between two strings is defined
as the minimum number of atomic operations (insertions, deletions, or
substitutions) required to transform one string into another.

\subsubsection{Edit Operations}\label{edit-operations}

Let \(X\) and \(Y\) be two strings. We consider three operations: 1.
\textbf{Insertion (I):} Add a character to the source string. 2.
\textbf{Deletion (D):} Remove a character from the source string. 3.
\textbf{Substitution (S):} Replace one character with another.

The total cost is the sum of these operations. Depending on the
application, specific operations may be weighted differently (e.g., in
some phonetic models, a substitution might count as 2, while an
insertion counts as 1).

\subsection{Dynamic Programming
Approach}\label{dynamic-programming-approach}

Computing the edit distance for every pair of strings in a large
vocabulary is computationally expensive. We use \textbf{Dynamic
Programming (DP)} to solve this by storing the results of subproblems in
a matrix \(D\) of size \((n+1) \times (m+1)\), where \(n\) and \(m\) are
the lengths of strings \(X\) and \(Y\).

\subsubsection{Formalization}\label{formalization}

Let \(D[i, j]\) denote the edit distance between the prefix
\(X[1 \dots i]\) and the prefix \(Y[1 \dots j]\).

The value \(D[i, j]\) is computed using the following recurrence
relation:

\[
D[i, j] = \min \begin{cases} 
D[i-1, j] + 1 & \text{(Deletion)} \\
D[i, j-1] + 1 & \text{(Insertion)} \\
D[i-1, j-1] + \text{cost}(X[i], Y[j]) & \text{(Substitution)}
\end{cases}
\]

Where the cost function for substitution is: *
\(\text{cost}(X[i], Y[j]) = 0\) if \(X[i] = Y[j]\) *
\(\text{cost}(X[i], Y[j]) = 1\) if \(X[i] \neq Y[j]\)

\subsubsection{Boundary Conditions}\label{boundary-conditions}

The base cases for the matrix are: * \(D[0, 0] = 0\) (Two empty strings
have distance 0) * \(D[i, 0] = i\) (Transforming a string of length
\(i\) to an empty string requires \(i\) deletions) * \(D[0, j] = j\)
(Transforming an empty string to a string of length \(j\) requires \(j\)
insertions)

\subsubsection{Computational Complexity}\label{computational-complexity}

\begin{itemize}
\tightlist
\item
  \textbf{Time Complexity:} \(O(n \times m)\), where \(n\) and \(m\) are
  the lengths of the two strings.
\item
  \textbf{Space Complexity:} \(O(n \times m)\) to store the matrix,
  though it can be optimized to \(O(\min(n, m))\) if only the final
  distance value is required.
\end{itemize}

\subsection{Application in NLP}\label{application-in-nlp}

\begin{itemize}
\tightlist
\item
  \textbf{Spell Checking:} Given a user input \(X\), finding the word
  \(W\) in the vocabulary \(V\) that minimizes \(D(X, W)\).
\item
  \textbf{Query Expansion:} Identifying similar terms for search
  optimization.
\item
  \textbf{Limitations:} Simple MED treats all character edits equally.
  In practice, some typos are more likely than others (e.g., hitting a
  key adjacent to the correct one on a QWERTY keyboard). Modern systems
  often use weighted edit distance matrices to account for keyboard
  topography or phonetic similarity (e.g., Soundex or Metaphone).
\end{itemize}

\subsubsection{Key Takeaways}\label{key-takeaways-4}

\begin{itemize}
\tightlist
\item
  \textbf{Minimum Edit Distance} provides a robust, deterministic way to
  quantify string dissimilarity.
\item
  \textbf{Dynamic Programming} is essential to reduce the complexity of
  the calculation from exponential to quadratic, making it feasible for
  real-time applications.
\item
  While subword tokenization provides an alternative to dictionary-based
  lookups, Edit Distance remains a foundational tool for correcting
  orthographic noise and identifying variants in low-resource settings.
\end{itemize}

\newpage

\section{N-grams and Sequential
Modeling}\label{n-grams-and-sequential-modeling}

In Natural Language Processing, an \textbf{N-gram} is a contiguous
sequence of \(N\) items from a given sample of text. These items can be
characters, syllables, or words, depending on the application. N-grams
are the building blocks of statistical language models, allowing us to
represent the structure of language based on local context.

\subsection{Definition and Taxonomy}\label{definition-and-taxonomy}

An N-gram model predicts the likelihood of the next element in a
sequence based on the preceding \(N-1\) elements. The size \(N\)
determines the scope of the ``memory'' of the model.

\subsubsection{Terminology}\label{terminology}

\begin{itemize}
\tightlist
\item
  \textbf{Unigram (\(N=1\)):} A single element (e.g., ``The''). This
  treats each word as an independent event, ignoring order.
\item
  \textbf{Bigram (\(N=2\)):} A sequence of two adjacent elements (e.g.,
  ``The cat''). This accounts for local pairwise dependencies.
\item
  \textbf{Trigram (\(N=3\)):} A sequence of three adjacent elements
  (e.g., ``The cat sat'').
\item
  \textbf{Quadrigram (\(N=4\)):} A sequence of four adjacent elements
  (e.g., ``The cat sat on'').
\end{itemize}

\subsection{Probabilistic Formulation}\label{probabilistic-formulation}

The objective of an N-gram model is to estimate the probability of a
word \(w_n\) given its history of length \(N-1\):

\[P(w_n | w_1, w_2, \dots, w_{n-1}) \approx P(w_n | w_{n-N+1}, \dots, w_{n-1})\]

\subsubsection{The Markov Assumption}\label{the-markov-assumption}

N-gram models operate under a \textbf{Markov Assumption}, which posits
that the probability of a future state depends only on a limited window
of recent history. * For a bigram model (\(N=2\)), the model assumes
\(P(w_n | w_{1:n-1}) \approx P(w_n | w_{n-1})\). * As \(N\) increases,
the model captures more long-range dependencies, but the number of
unique N-gram types (and thus the memory and training data requirements)
grows exponentially.

\subsection{Practical Considerations}\label{practical-considerations}

\subsubsection{Sparsity}\label{sparsity}

As \(N\) increases, the number of possible N-gram combinations \(|V|^N\)
grows rapidly. This leads to \textbf{data sparsity}, where many valid
sequences in a language never appear in the training corpus (assigned a
probability of zero). To mitigate this, practitioners use: *
\textbf{Smoothing:} Assigning a small probability mass to unseen N-grams
(e.g., Laplace smoothing, Good-Turing estimation). * \textbf{Backoff:}
Using a lower-order N-gram (e.g., a bigram) if the higher-order N-gram
(e.g., a trigram) has not been observed. * \textbf{Interpolation:}
Mixing probabilities from different N-gram models (e.g.,
\(0.6 \times P_{trigram} + 0.3 \times P_{bigram} + 0.1 \times P_{unigram}\)).

\subsubsection{N-grams vs.~Subword
Models}\label{n-grams-vs.-subword-models}

While N-grams operate on discrete word units, subword tokenization (like
BPE) acts as a specialized form of N-gram induction where tokens are
learned based on character-sequence frequency rather than pre-defined
word boundaries.

\subsubsection{Key Takeaways}\label{key-takeaways-5}

\begin{itemize}
\tightlist
\item
  \textbf{Contiguity:} N-grams capture sequential local dependencies by
  definition.
\item
  \textbf{Trade-off:} Small \(N\) values are robust but capture minimal
  context; large \(N\) values provide more linguistic nuance but suffer
  severely from data sparsity.
\item
  \textbf{Language Modeling:} N-grams provide the foundational framework
  for estimating the probability of word sequences, a core task in
  speech recognition, machine translation, and text generation.
\item
  \textbf{Scalability:} The choice of \(N\) must be balanced against the
  available corpus size to avoid the ``curse of dimensionality''
  inherent in counting sequences.
\end{itemize}

\newpage

\section{Language Models: Foundations and
Probability}\label{language-models-foundations-and-probability}

A Language Model (LM) is a computational system that assigns a
probability distribution over sequences of words or tokens. By learning
patterns from large-scale text data, an LM allows us to quantify the
likelihood of a given string, effectively capturing the structure and
``naturalness'' of human language.

\subsection{Formalization and
Probability}\label{formalization-and-probability}

The fundamental goal of a language model is to compute the probability
of a sequence of \(N\) words \(W = (w_1, w_2, \dots, w_n)\). Using the
\textbf{Chain Rule of Probability}, we decompose this joint probability
into a product of conditional probabilities:

\[P(w_1, w_2, \dots, w_n) = \prod_{i=1}^{n} P(w_i | w_1, w_2, \dots, w_{i-1})\]

In this formulation, \(P(w_i | w_1, \dots, w_{i-1})\) represents the
probability of word \(w_i\) occurring given all preceding words (the
``history'').

\subsubsection{The N-gram Approximation}\label{the-n-gram-approximation}

Since computing \(P(w_i | w_1, \dots, w_{i-1})\) exactly is intractable
due to the infinite variety of language and the sparsity of long
histories, we apply the \textbf{Markov Assumption}. We approximate the
probability by considering only the most recent \(N-1\) words:

\[P(w_i | w_1, \dots, w_{i-1}) \approx P(w_i | w_{i-N+1}, \dots, w_{i-1})\]

For example, in a \textbf{Bigram model} (\(N=2\)):
\[P(w_i | w_1, \dots, w_{i-1}) \approx P(w_i | w_{i-1})\]

\subsection{Estimation and Training}\label{estimation-and-training}

To estimate these probabilities, we use the Maximum Likelihood
Estimation (MLE) approach based on observed frequencies in a corpus:

\[P(w_i | w_{i-N+1:i-1}) = \frac{C(w_{i-N+1:i-1}, w_i)}{C(w_{i-N+1:i-1})}\]

Where: * \(C(\cdot)\) denotes the count of the sequence in the corpus. *
The denominator is the total count of the history sequence.

\subsubsection{Challenges in Estimation}\label{challenges-in-estimation}

\begin{enumerate}
\def\labelenumi{\arabic{enumi}.}
\tightlist
\item
  \textbf{Corpus Dependency:} The quality and domain of the training
  corpus dictate the model's performance. A model trained on legal
  documents will assign low probability to casual conversation.
\item
  \textbf{Infinite Language:} Language is creative and open-ended. Even
  with massive corpora, many valid sequences are never observed (Zero
  Probability Problem), leading to the need for smoothing techniques.
\item
  \textbf{Data Sparsity:} As \(N\) increases, the number of possible
  unique sequences grows exponentially (\(|V|^N\)), meaning most N-grams
  will have zero counts in any finite dataset.
\end{enumerate}

\subsection{Applications of Language
Models}\label{applications-of-language-models}

Language models serve as the backbone for various NLP tasks: *
\textbf{Speech Recognition:} Disambiguating acoustic signals by
assigning high probability to likely word sequences. * \textbf{Machine
Translation:} Ranking candidate translations based on their fluency in
the target language. * \textbf{Spelling Correction:} Selecting the most
probable intended word among candidates based on context. *
\textbf{Part-of-Speech Tagging:} Determining the correct syntactic tag
for a word based on its sequence of neighbors.

\subsubsection{Key Takeaways}\label{key-takeaways-6}

\begin{itemize}
\tightlist
\item
  \textbf{Language Models} define a probability distribution over
  strings, allowing for generative and interpretive capabilities.
\item
  The \textbf{Chain Rule} allows for the decomposition of sequence
  probability, while the \textbf{Markov Assumption} makes the
  calculation computationally feasible through N-gram approximations.
\item
  \textbf{Probability vs.~Fluency:} An ideal LM assigns high probability
  to grammatically correct and semantically logical sequences, and very
  low probability to nonsensical or rare sequences.
\item
  Estimation relies on frequency counting, but requires advanced
  techniques like smoothing to handle the inherent sparsity of human
  language.
\end{itemize}

\newpage

\section{Statistical Language Modeling and
Estimation}\label{statistical-language-modeling-and-estimation}

\subsection{Relative Frequency Estimation (Maximum
Likelihood)}\label{relative-frequency-estimation-maximum-likelihood}

To derive a probability distribution for an N-gram model, we rely on the
\textbf{Relative Frequency Estimate}. This approach, known as
\textbf{Maximum Likelihood Estimation (MLE)}, chooses the parameter
values that maximize the likelihood of the observed training data.

\subsubsection{The MLE Formulation}\label{the-mle-formulation}

Given a corpus, we normalize the counts of N-gram occurrences by the
counts of their prefixes.

For a bigram model (\(N=2\)), the probability of \(w_i\) given
\(w_{i-1}\) is:
\[P(w_i | w_{i-1}) = \frac{C(w_{i-1}, w_i)}{C(w_{i-1})}\]

For the general case (\(N\)-gram), the probability of \(w_i\) given a
context of length \(N-1\) is:
\[P(w_i | w_{i-N+1}, \dots, w_{i-1}) = \frac{C(w_{i-N+1}, \dots, w_{i-1}, w_i)}{C(w_{i-N+1}, \dots, w_{i-1})}\]

\subsubsection{Sentence Boundaries}\label{sentence-boundaries}

To allow the model to learn the probability of a sentence starting or
ending, we augment the text by bracketing each sentence with special
start (\texttt{\textless{}s\textgreater{}}) and end
(\texttt{\textless{}/s\textgreater{}}) tokens. This ensures that the
model can learn the distribution of valid sentence openers and closers,
treating them as part of the sequence.

\subsection{Statistical Model Training and Data
Splitting}\label{statistical-model-training-and-data-splitting}

Statistical models like N-grams require careful handling of data to
ensure that the learned parameters generalize to unseen text rather than
simply memorizing the training material.

\subsubsection{The Three-Way Data Split}\label{the-three-way-data-split}

To evaluate a model's performance accurately, we partition the available
data into three distinct, non-overlapping sets:

\begin{enumerate}
\def\labelenumi{\arabic{enumi}.}
\tightlist
\item
  \textbf{Training Set (80\%):} The data used to calculate the counts
  \(C(\cdot)\) and populate the N-gram probability tables. The model
  ``learns'' from this data.
\item
  \textbf{Development (Dev) Set (10\%):} Used for tuning
  hyperparameters---such as the choice of \(N\) in an N-gram model, or
  the selection of smoothing parameters. This set helps prevent
  overfitting to the training data.
\item
  \textbf{Test Set (10\%):} A ``held-out'' dataset containing fresh data
  that the model has never encountered. The performance on this set
  provides an unbiased estimate of the model's generalization
  capabilities.
\end{enumerate}

\subsubsection{Key Considerations}\label{key-considerations}

\begin{itemize}
\tightlist
\item
  \textbf{Non-Intersection:} It is strictly required that the test set
  shares no data with the training or development sets. Evaluation on
  training data (``cheating'') leads to overly optimistic performance
  reports.
\item
  \textbf{Generalization:} The fundamental goal is to build a model that
  captures the underlying linguistic patterns, enabling it to process
  novel sequences accurately.
\end{itemize}

\subsubsection{Key Takeaways}\label{key-takeaways-7}

\begin{itemize}
\tightlist
\item
  \textbf{MLE} is the standard approach for estimating N-gram
  probabilities by normalizing counts based on observed occurrences in
  the training corpus.
\item
  \textbf{Sentence delimiters (\texttt{\textless{}s\textgreater{}},
  \texttt{\textless{}/s\textgreater{}})} are essential for modeling the
  probability of entire sequences rather than just word transitions.
\item
  \textbf{Data Partitioning} into Training, Development, and Test sets
  is mandatory to develop models that generalize to unseen, real-world
  data and to tune parameters without overfitting.
\end{itemize}

\newpage

\section{Perplexity and Evaluation of Language
Models}\label{perplexity-and-evaluation-of-language-models}

Once a language model is trained, it is essential to quantify its
performance. A primary metric for this, which reflects how well a model
fits a test dataset, is \textbf{Perplexity (PP)}.

\subsection{Understanding Perplexity}\label{understanding-perplexity}

Perplexity measures how ``surprised'' a language model is by a test set.
A model that assigns a higher probability to the sentences in the test
set is considered a better model, and thus, it will have lower
perplexity.

Formally, perplexity is the inverse probability of the test set,
normalized by the number of words \(n\).

\subsubsection{Formal Definition}\label{formal-definition}

For a sequence of words \(W = w_1, w_2, \dots, w_n\), the perplexity is
defined as: \[PP(W) = P(w_1, \dots, w_n)^{-1/n}\]

Using the Chain Rule of probability, this can be expanded as:
\[PP(W) = \left( \prod_{i=1}^{n} P(w_i | w_1, \dots, w_{i-1}) \right)^{-1/n}\]

\subsubsection{Perplexity for a Bigram
Model}\label{perplexity-for-a-bigram-model}

In a bigram model, the conditional probability
\(P(w_i | w_1, \dots, w_{i-1})\) simplifies to \(P(w_i | w_{i-1})\).
Consequently, the perplexity becomes:
\[PP(W) = \sqrt[n]{\frac{1}{\prod_{i=1}^{n} P(w_i | w_{i-1})}}\]

\begin{itemize}
\tightlist
\item
  \textbf{Interpretation:} Perplexity is the \textbf{geometric mean of
  the inverse probabilities}.
\item
  \textbf{Intuition:} A model with a perplexity of \(k\) is effectively
  as ``confused'' as if it had to choose uniformly at random between
  \(k\) possible next words at each position in the test set. Therefore,
  lower perplexity is always better.
\end{itemize}

\subsection{Qualitative Evaluation}\label{qualitative-evaluation}

Beyond quantitative metrics like perplexity, N-gram models are
frequently evaluated via \textbf{qualitative inspection}:

\begin{enumerate}
\def\labelenumi{\arabic{enumi}.}
\tightlist
\item
  \textbf{Sentence Generation:} We can use the trained model to sample
  sequences of words. By observing the generated output, we can
  qualitatively assess the fluency and coherence of the model.
\item
  \textbf{Order Sensitivity:} The quality of generated text varies
  significantly with the choice of \(N\):

  \begin{itemize}
  \tightlist
  \item
    \textbf{Unigrams (\(N=1\)):} Generate ``word soup'' because they
    ignore the order and relationships between words.
  \item
    \textbf{Bigrams (\(N=2\)):} Capture simple local dependencies; the
    output often makes local sense but lacks long-term coherence.
  \item
    \textbf{Trigrams/Quadrigrams (\(N=3, 4\)):} Produce increasingly
    readable text as they capture larger, more contextually relevant
    chunks of language, though they are more prone to verbatim
    memorization of the training set if \(N\) is too large.
  \end{itemize}
\end{enumerate}

\subsubsection{Key Takeaways}\label{key-takeaways-8}

\begin{itemize}
\tightlist
\item
  \textbf{Perplexity} is a crucial metric that evaluates a model's fit
  on unseen data; lower scores indicate higher predictive accuracy.
\item
  \textbf{Geometric Interpretation:} Perplexity effectively measures the
  branching factor of the model's predictions.
\item
  \textbf{Qualitative vs.~Quantitative:} While perplexity provides a
  scalar score, generating text serves as a sanity check to verify if
  the model has learned the structural and syntactic nuances of the
  target language.
\end{itemize}

\newpage

\section{Smoothing in Statistical Language
Models}\label{smoothing-in-statistical-language-models}

Maximum Likelihood Estimation (MLE) assigns probabilities based strictly
on observed frequencies in a training corpus. Because language is
open-ended and the training data is necessarily finite, many valid
N-grams (sequences that could exist in a language) will never be
observed in the training set.

This leads to the \textbf{Zero Probability Problem}: if a test sequence
contains a single unseen N-gram, the entire sequence probability becomes
zero. \textbf{Smoothing} (or discounting) is the collection of
techniques used to reallocate a small portion of the probability mass
from observed N-grams to unseen N-grams.

\subsection{Smoothing Techniques}\label{smoothing-techniques}

\subsubsection{1. Laplace (Add-One)
Smoothing}\label{laplace-add-one-smoothing}

The simplest approach is to add a constant count (usually 1) to all
possible N-grams in the vocabulary, even if their observed count is
zero.

For a bigram model:
\[P_{Laplace}(w_i | w_{i-1}) = \frac{C(w_{i-1}, w_i) + 1}{C(w_{i-1}) + |V|}\]

\begin{itemize}
\tightlist
\item
  \textbf{Mechanism:} Adding 1 to the numerator of all counts and
  \(|V|\) to the denominator ensures that the sum of all probabilities
  remains 1.
\item
  \textbf{Limitation:} It is generally too aggressive; it steals too
  much probability mass from frequent words and assigns it to rare ones,
  which often leads to poor performance.
\item
  \textbf{Add-Lambda (\(\lambda\)):} A refinement where one adds a small
  fraction \(\lambda\) (e.g., 0.1) instead of 1, providing finer control
  over the smoothing.
\end{itemize}

\subsubsection{2. Good-Turing Smoothing}\label{good-turing-smoothing}

Good-Turing estimation uses the count of things that have occurred once
(\(N_1\)) to estimate the probability of things that have occurred zero
times (\(N_0\)).

\begin{itemize}
\tightlist
\item
  \textbf{Logic:} If \(N_r\) is the number of N-grams that appear \(r\)
  times, we estimate the count of N-grams that appear \(r\) times using
  the count of N-grams that appear \(r+1\) times.
\item
  \textbf{Application:} It is particularly effective for estimating the
  probability mass for unseen items.
\end{itemize}

\subsubsection{3. Interpolation}\label{interpolation}

Interpolation combines higher-order N-gram models with lower-order
models. It assumes that if a trigram is unseen, we can fall back on the
bigram or unigram probability as a proxy.

\[P_{interp}(w_n | w_{n-2}, w_{n-1}) = \lambda_1 P(w_n | w_{n-2}, w_{n-1}) + \lambda_2 P(w_n | w_{n-1}) + \lambda_3 P(w_n)\]

Where \(\sum \lambda_i = 1\). This allows the model to leverage the
specificity of higher orders when available, and the robustness of lower
orders when data is sparse.

\subsubsection{4. Backoff Strategies (Katz
Backoff)}\label{backoff-strategies-katz-backoff}

Unlike interpolation, which always uses a mix of models,
\textbf{Backoff} uses the most specific model available.

\begin{itemize}
\tightlist
\item
  If \(C(w_{n-2}, w_{n-1}, w_n) > 0\), use the trigram probability.
\item
  Otherwise, ``back off'' to the bigram probability \(P(w_n | w_{n-1})\)
  multiplied by a scaling factor to ensure the total distribution
  remains normalized.
\item
  If the bigram is also unseen, back off to the unigram probability.
\end{itemize}

\subsubsection{Key Takeaways}\label{key-takeaways-9}

\begin{itemize}
\tightlist
\item
  \textbf{Data Sparsity:} Smoothing is a necessary correction for the
  inherent limitations of empirical frequency estimation.
\item
  \textbf{Trade-off:} Effective smoothing balances the need to assign
  non-zero probability to unseen events without sacrificing the
  predictive power of high-frequency, well-observed sequences.
\item
  \textbf{Evolution:} While these ``pre-neural'' techniques were
  essential for classical N-gram models, they paved the way for modern
  methods like embedding-based smoothing, where similarity in vector
  space allows for the generalization to unseen contexts.
\end{itemize}

\end{document}
